\appendix
\section{Measurability and translation averages}\label{sec:background}
We give the measure-theoretic details used in the component argument.
They require no moment bound.

\begin{lemma}\label{lem:standard-borel}
The space \(\mathcal X=\mathcal S'(\R^3;V)\), with its evaluation
sigma-algebra, is standard Borel. Translation is a jointly measurable
action on this space.
\end{lemma}
\begin{proof}
Choose a countable rational vector space \(\mathcal D_0\) dense in
\(\mathcal S(\R^3;V)\), and increasing norms \(q_n\) defining the
Schwartz topology. For \(n\ge1\), put
\[
 X_n=\{F:|F(f)|\le n q_n(f)\text{ for every }f\in\mathcal S\}.
\]
Continuity of each distribution gives \(\mathcal X=\bigcup_nX_n\).
The evaluation map
\(\iota(F)=(F(f))_{f\in\mathcal D_0}\) is injective into the
Polish space \(\R^{\mathcal D_0}\).
Its image on \(X_n\) is exactly the set of sequences satisfying
\[
 z_{sf+tg}=sz_f+tz_g,\qquad |z_f|\le n q_n(f)
 \quad(f,g\in\mathcal D_0,\ s,t\in\Q).
\]
These are closed conditions on countably many coordinates.
Conversely, the bound extends such a rational linear functional
continuously and uniquely to the Schwartz space; density and
continuity make the extension real linear.
Thus \(\iota(\mathcal X)\) is a countable union of closed sets.
Every evaluation is measurable with respect to the
\(\mathcal D_0\) coordinates by approximation and samplewise
continuity. Hence \(\iota\) is a measurable isomorphism onto a
Borel subset of a Polish space, which is standard Borel
\cite[Theorem~2.1.22]{Durrett2019}.

For fixed \(f\), the function
\((a,F)\mapsto F(\tau_{-a}\cdot f)\) is continuous in \(a\) and
measurable in \(F\). Replace \(a\) by its coordinatewise lower
dyadic approximation \(2^{-n}\lfloor2^na\rfloor\).
The resulting functions are measurable on the product space and
converge pointwise to the original one. Since evaluations generate
the sigma-algebra, the translation action is jointly measurable.
\end{proof}

\begin{lemma}\label{lem:averages-translations}
Let \(\lambda\) be translation invariant. If \(J\subset\{1,2,3\}\)
is nonempty, the averages
\[
 M_T^J A=T^{-|J|}\int_{[0,T]^J}U(a)A\,\dd a
\]
converge in \(L^2(\lambda)\) to the orthogonal projection \(P_J\)
onto the observables fixed by the translations in the coordinate
subspace spanned by \(\{e_j:j\in J\}\).
This projection is conditional expectation onto the corresponding
invariant sigma-algebra. All \(M_T^J\) are contractions on
\(L^p(\lambda)\), \(1\le p<\infty\).
\end{lemma}
\begin{proof}
Joint measurability and Fubini give the averages, and Jensen's
inequality gives the \(L^p\) contraction. The representation \(U\)
is strongly continuous by bounded convergence on local characters
and their density in \(L^2(\lambda)\).
Let \(\mathcal K\) be the closed span of
\(B-U(b)B\), where \(b\) is in the coordinate subspace.
Its orthogonal complement is the invariant subspace.
The averages fix that subspace. The symmetric-difference bound
for a cube and its translate gives
\[
 \|M_T^J(B-U(b)B)\|_2
 \le\frac{2|b|_1}{T}\|B\|_2.
\]
Thus the averages tend to zero on a dense subspace of
\(\mathcal K\), and their uniform norm bound extends this
convergence to all of \(\mathcal K\).
This proves the asserted orthogonal-projection limit.

The invariant subspace equals the \(L^2\) space of the invariant
sigma-algebra: indicators of invariant events are fixed, and
the inverse images of Borel sets under a fixed observable are
invariant modulo null sets. Therefore its orthogonal projection
is the stated conditional expectation. The coordinate projections
commute, since their averaging operators commute.
\end{proof}
